% Third part of holding_prm_report.tex (input from there): the roadmap as a
% backend of plan_constrained_dual_arm, used in scenes that change from bar to bar.
% Tables come from docs/figures/make_prm_figures.py (robarch_* functions).

% =============================================================================
\section{A precomputed roadmap behind \texttt{plan\_constrained\_dual\_arm}, in a changing scene}
\label{sec:scene}
% =============================================================================

Sections~\ref{sec:build}--\ref{sec:refine} built and refined roadmaps inside the
benchmark's scenes, which hold only the robot, its tools and the carried bar. On
site the scene is different for every bar: the bars built so far, their joints,
and sometimes a support robot stand around the goal, and the base parks a few
centimetres off its planned pose. This section puts the roadmap behind the
production entry point and makes it work there. The idea is the one from the
meeting notes: \emph{precompute} a roadmap per bar that knows only the robot, its
tools and the bar; \emph{online}, place it at the real base pose and check only
what the precompute could not know.

\subsection{What is reused, and what has to be redone}

\begin{center}
\footnotesize
\begin{tabular}{@{}p{4.6cm} p{5.0cm} p{5.6cm}@{}}
\toprule
 & depends on & online \\
\midrule
roadmap nodes and edges & the bar, its two grasps & reused as they are \\
collisions among robot, tools, bar and the joint halves on the bar & the same & reused (verdicts of the precompute) \\
exact waypoints of each edge (tracking) & the same & reused \\
the goal as seen from the robot & base pose, goal in the world & solved fresh: a few IK calls \\
collisions with built bars, joints, support robots & base pose, the scene & checked fresh, only on what a route uses \\
\bottomrule
\end{tabular}
\end{center}

\noindent
A node is twelve joint values, so it does not depend on where the base stands;
neither does anything that rides on the arms. The roadmap is therefore built once,
\emph{in the robot's frame}, and ``moving it to the real base pose'' needs no
computation on the roadmap at all: the planner reads the two arm bases from the
live scene and expresses the world goal with them,
\[
  {}^{\text{base}}T_{\text{bar}}^{\text{goal}} \;=\;
  \bigl({}^{\text{world}}T_{\text{base}}\bigr)^{-1}\; {}^{\text{world}}T_{\text{bar}}^{\text{goal}} .
\]
A base parked elsewhere moves the goal as seen from the robot, and moves the
scene; nothing else. There is one roadmap per bar, never several to stitch.

\subsection{Splitting the collision check}
\label{sec:split}

The production check (\code{planner.check_collision} of compas\_fab) tests
everything against everything in one call and has no per-body filter. The split
is made by editing a private copy of the cell state, so compas\_fab itself is
unchanged (\code{sampling.make_collision_fns}):

\begin{description}[nosep,leftmargin=1.5em]
  \item[full] the production check, as before.
  \item[robot side] every body and tool that does not ride on the robot is hidden.
        What is left is exactly what the precompute scene tests.
  \item[scene only] every carried body is allowed to touch all robot links, all
        other carried bodies and the arm tools; the arm tools are allowed to touch
        all robot links; and the robot's self-collision pass is skipped. What is
        left is robot links, tools and carried bodies against the bodies and tools
        standing in the scene. Allowances the design authored toward the scene
        (the bar may touch the joint halves it will sit in) are kept.
\end{description}

\noindent
Two facts make the reuse sound, and both are tested on a real scene (the last
test of \path{test_holding_prm.py}): the full check says
``collides'' exactly when the robot-side check or the scene-only check does; and
the precompute scene, with the base at the origin and nothing built, gives the
same verdict as the robot-side check in the real scene at the real base pose.
Table~\ref{tab:robarchsplit} gives the result and the time of each check.

\paragraph{A signature guards the reuse.} The robot-side verdicts are valid only
for the same load. The roadmap stores a signature of what they depend on (the
bar, every carried body with its link, attachment frame and allowed contacts
among the carried bodies, the arm tools' allowed links); the planner compares it
with the live start state and the caller's right-hand grasp, and stops with
\code{RoadmapMismatch}, listing the differences, when they do not match. A
roadmap built from a benchmark library has no signature and is refused. There
is no fallback to another roadmap or to the BiRRT.

\subsection{Precompute for one exported bar}

\code{prm_pilot.py build --problem-dir <export> --bar B4} builds the roadmap of
one bar from the exported action files, in the trimmed scene of
\code{benchmark build-cells} (robot, tools, the bar and its joint halves).

\begin{description}[nosep,leftmargin=1.5em]
  \item[goal anchors] the planned goal pose of the carry movement
        (\code{target_ee_frames} of the jointing file, seen from the planned base
        pose), plus the same world goal seen from 12 bases shifted by up to
        10\,cm and turned by up to $5^\circ$ (\code{--base-error-m},
        \code{--base-error-deg}, \code{--base-samples}). Each pose is expanded to
        all its collision-free branch pairs. The roadmap is then dense around a
        goal \emph{region}, so that the real goal, which will coincide with none of
        the 13, has neighbours.
  \item[start anchors] the 18 carry poses of the family found in
        Section~\ref{sec:e5} (\path{prm/carry/carry_v1/starts.json}), every
        branch pair.
  \item[the rest] sampling, edges and exact tracking of every edge within 15\,mm
        (\code{prm_pilot.py track}), as in Sections~\ref{sec:build}
        and~\ref{sec:refine}.
\end{description}

\subsection{The online plan}
\label{sec:online}

\begin{algo}{\textsc{PlanOnRoadmap} --- \texttt{online.plan\_on\_roadmap}}{alg:online}
\Function{PlanOnRoadmap}{live scene and start state, world goal pose of the bar, roadmap folder}
  \State load the roadmap; \textbf{stop} unless its signature matches the load in the scene \Comment{no fallback}
  \State read the arm bases from the live scene; check numpy FK against the simulator ($10^{-5}$\,m)
  \State workers: this process (1 worker), or $n$ processes that open the scene from files
         \Comment{cell cached by geometry}
  \State $G \gets$ every collision-free branch pair at the goal pose \Comment{full check, in the scene}
  \State $S \gets$ the roadmap's carry-pose nodes that pass the scene-only check
  \State join new goal nodes to their nearest roadmap nodes; load the night's tracking verdicts
  \For{tolerance $t$ in 10, 20, 30\,mm}
    \Loop
      \State route $\gets$ shortest route $S \to G$ over edges within $t$ not known bad; \textbf{break} if none
      \State edges tracked overnight: scene-only check of their stored waypoints \Comment{parallel, lazy}
      \State other edges (they touch a new goal node): track now, full check \Comment{parallel}
      \State drop the edges that failed; \textbf{continue} unless the whole route is usable
      \State join the edges' waypoints into one exact path \Comment{Section~\ref{sec:refine}}
      \State walk it at the caller's step (2\,mm) and full-check every fine waypoint \Comment{parallel}
      \State \Return the fine path if clean; else drop the edge that failed and \textbf{continue}
    \EndLoop
  \EndFor
  \State one tube round around the best route left (Section~\ref{sec:refine}), then the loop once more
  \State \Return failure, with the smallest tolerance a remaining route would need
\EndFunction
\end{algo}

\noindent
The scene check is \emph{lazy}: only edges on a candidate route are checked, and
only against the scene, because their robot side was settled overnight. An edge
keeps a separate chain of waypoints for each direction; a direction that touches
the scene is dropped on its own. The final check replaces the production repair
of Section~\ref{sec:refine} (which ran in one process and took 30--190\,s per
path): the same fine walk and the same collision check, spread over the workers,
and a failure costs one edge and another search round rather than the plan.

\paragraph{A new base pose, or a newly built bar.} Nothing in the roadmap is
redone. The goal nodes are solved again (the goal moved as seen from the robot)
and joined to the roadmap; scene verdicts are asked again, lazily. The cached
scene files are keyed by the scene's geometry, so a new base pose rewrites only
the small state file, not the cell.

\subsection{Switching planners}

\code{plan_constrained_dual_arm} keeps its signature and gains two keywords:
\code{backend="birrt"} (default, unchanged) or \code{"prm"}, and
\code{prm_options} (\code{roadmap_dir}, \code{workers}, \code{tolerances_mm},
\code{tube_rounds}). The roadmap branch returns the same \code{(path, info)}; an
unknown backend, a missing roadmap or a roadmap for another load raise. After
the roadmap plan the path goes through the same optional repair and the same
shortcut smoothing as a BiRRT path.

\code{scripts/headless_bar_action_planner.py} plans the bar carry of an exported
jointing file with either planner:
\begin{quote}\footnotesize
\code{headless_bar_action_planner.py <data> --problem 260920_RobArch_demo_revamp}\\
\hspace*{1.5em}\code{--bar-action B4__J.json --movement cdfm --cdfm-backend prm --prm-workers 8}
\end{quote}
(\code{--movement cdfm} finds the carry by its type and reads home from the
release file; \code{--prm-run} names another roadmap than
\code{<problem>__<bar>}; \code{--save-tag} and \code{--save-dir} keep results of
different runs apart; \code{--base-shift} plans as if the base parked off its
planned pose.) \code{scripts/prm_check_solved.py} then checks a saved path on its
own: production repair at 2\,mm with the calibrated solver, the $10^\circ$ gate
and the full collision check in the scene of the saved start state.

\paragraph{Looking at a plan.} Two ways. \code{--replay-file <saved motion file> --gui}
replays any saved path in the planner script's PyBullet window, in the scene of the
saved start state, with a slider along the path. And every online plan is saved as a
query of its roadmap (\path{prm/<run>/queries/<bar>_<tag>.json.gz}, with the scene
it was planned in): after \code{prm_pilot.py export --name <run>}, which also writes
the robot's and the scene's 3D files, the roadmap page
(\path{http://localhost:8000/prm/<run>/dashboard/} over a web server on the data root)
draws the built bars, joint halves and support robots around the robot, marks the
edges the scene blocked, and walks the plan's exact path with its stepper; the
holding error it measures along that path is zero. Under the stepper the page
plots, for the exact path or for a node's best route alike, the holding error at
the commanded points (one dot per point, so the density of the data is visible),
the holding error along a dense joint-space interpolation of the same path
($0.1^\circ$ per joint, the step of the planner's own path checker), which
approximates what the robot shows while its controller tracks the path between
commanded points, and the twelve joint angles along the path. On the exact B4 path
(1525 points, at most $0.3^\circ$ between neighbours) the dense interpolation adds
nothing measurable; on a node route interpolated at the roadmap's $2.5^\circ$ it
rises by a few hundredths of a millimetre over the value at the points.

\subsection{Test bed: B1, B3 and B4 of the RobArch demo export}
\label{sec:scenetest}

The three bars are the first three the two-arm robot carries in
\path{data_design_study/260920_RobArch_demo_revamp}; their scenes grow as the
structure does (Table~\ref{tab:robarchpre}). The movement planned is the bar carry
embedded in each jointing file (movement \code{J_M3}, the constrained two-arm
free movement), from a carry pose the planner
chooses to the authored goal pose, any arm posture.

\begin{table}[h]
\centering
\footnotesize
\setlength{\tabcolsep}{4pt}
\caption{The three bars, their scenes at carry time, and the roadmap precomputed for
each (6000 sampled nodes plus anchors; 8 workers). B3's and B4's tracking ran while
other jobs used the machine; B1's times are from an idle machine.}
\label{tab:robarchpre}
\IfFileExists{figures/prm_table_robarch_precompute.tex}{\begin{tabular}{@{}>{\raggedright\arraybackslash}p{7.4cm} rrr@{}}
\toprule
bar & B1 & B3 & B4 \\
\midrule
bar length (m) & 1.80 & 2.00 & 0.70 \\
bodies riding on the arms: bar and joint halves (count) & 9 & 8 & 3 \\
scene at carry time: built bars and joint halves (count) & 0 & 9 & 17 \\
scene at carry time: support robots (count) & 2 & 2 & 2 \\
\midrule
goal poses the arms can hold: planned + 12 shifted bases & 11 of 13 & 13 of 13 & 13 of 13 \\
carry poses the arms can hold & 10 of 18 & 18 of 18 & 17 of 18 \\
nodes (count) & 6352 & 6460 & 6342 \\
collision-free edges (count) & 35461 & 34668 & 38689 \\
components of the free graph (count) & 243 & 302 & 224 \\
goal nodes in the largest component & 52 of 154 & 55 of 206 & 30 of 96 \\
carry-pose nodes in the largest component & 54 of 197 & 101 of 253 & 112 of 245 \\
\midrule
build: sampling and edges (s) & 641 & 631 & 490 \\
build: workers & 8 & 8 & 8 \\
edges within 15 mm, tracked exactly (count) & 4094 & 8260 & 12849 \\
\quad of which track & 3720 & 7930 & 12567 \\
tracking (s) & 1556 & 3691 & 6860 \\
\bottomrule
\end{tabular}
}{(missing)}
\end{table}

\begin{table}[h]
\centering
\footnotesize
\setlength{\tabcolsep}{4pt}
\caption{The roadmap planner on the bar carry of each bar, in its real scene, with 8
workers (\code{--cdfm-backend prm}). Times are wall-clock seconds on an otherwise
idle machine; ``worker-seconds'' add up the time inside all workers.}
\label{tab:robarchplans}
\IfFileExists{figures/prm_table_robarch_plans.tex}{\begin{tabular}{@{}>{\raggedright\arraybackslash}p{7.4cm} rrr@{}}
\toprule
bar & B1 & B3 & B4 \\
\midrule
result: tolerance of the route used & 20 mm & 20 mm & 10 mm \\
carry-pose nodes free in the scene & 197 of 197 & 248 of 253 & 245 of 245 \\
carry pose the route starts from & horizontal\_1 & horizontal\_4 & vertical\_4 \\
nodes on the route; waypoints returned (count) & 7; 824 & 7; 760 & 10; 1323 \\
edges looked at, of which usable (count) & 6, 6 & 6, 6 & 23, 14 \\
\quad blocked by the scene (count) & 0 & 0 & 9 \\
\quad dropped in the final check (count) & 0 & 0 & 0 \\
tube rounds needed (count) & 0 & 0 & 0 \\
scene-only checks: count, worker-seconds & 373, 3.1 & 349, 4.7 & 1187, 17.7 \\
full checks: count, worker-seconds & 998, 26.3 & 1478, 45.8 & 1421, 41.6 \\
\midrule
start the workers (s) & 15.4 & 17.0 & 15.2 \\
read the roadmap; write or find the scene files (s) & 0.7 & 0.8 & 0.7 \\
goal postures and carry-pose nodes in the scene (s) & 1.5 & 2.9 & 3.6 \\
search: scene checks and tracking of route edges (s) & 6.5 & 6.9 & 3.2 \\
final check at 2 mm (s) & 4.2 & 8.4 & 7.3 \\
roadmap plan in all (s) & 28.4 & 36.0 & 30.1 \\
smoothing afterwards, in the caller's process (s) & 57.4 & 44.5 & 30.1 \\
whole planning call (s) & 88.9 & 83.6 & 63.3 \\
\midrule
independent check of the saved path & pass & pass & pass \\
path length: bar travel (m); bar turning (deg) & 1.44; 70 & 1.13; 109 & 1.17; 158 \\
path length: joint travel (deg) & 140 & 146 & 190 \\
lowest point of bar, joints and tools on the path (m) & -0.003 & 0.048 & 0.071 \\
\quad the same at the goal itself (m) & -0.003 & 0.048 & 0.071 \\
\bottomrule
\end{tabular}
}{(missing)}
\end{table}

\begin{table}[h]
\centering
\footnotesize
\setlength{\tabcolsep}{4pt}
\caption{The BiRRT on the same three movements (three runs each; a run may use up to
five attempts of 300\,s), the roadmap planner with 8 workers and with one, and the
roadmap planner without the night's tracking file. ``fail (t)'': no path after
$t$ seconds.}
\label{tab:robarchcompare}
\IfFileExists{figures/prm_table_robarch_compare.tex}{\begin{tabular}{@{}>{\raggedright\arraybackslash}p{7.4cm} rrr@{}}
\toprule
bar & B1 & B3 & B4 \\
\midrule
BiRRT: runs solved (of 3) & 3 & 3 & 3 \\
BiRRT: whole planning call of each run (s) & 183, 188, 185 & 64, 65, 63 & 1624, 504, 397 \\
BiRRT: solved paths that pass the independent check & 3 of 3 & 3 of 3 & 3 of 3 \\
BiRRT: lowest point on its paths (m) & -0.003 & 0.048 & 0.071 \\
BiRRT: bar travel of each path (m) & 0.41, 0.51, 0.41 & 1.15, 1.15, 1.15 & 1.49, 2.06, 2.51 \\
BiRRT: joint travel of each path (deg) & 128, 217, 126 & 144, 144, 144 & 570, 563, 779 \\
\midrule
roadmap, 8 workers: whole planning call (s) & 89 & 84 & 63 \\
\quad of which the roadmap plan (s) & 28 & 36 & 30 \\
\quad of which starting the workers (s) & 15 & 17 & 15 \\
roadmap, 1 worker: whole planning call (s) & 78 & 116 & 118 \\
\quad of which the roadmap plan (s) & 57 & 86 & 93 \\
\midrule
roadmap, 8 workers: search with the night's tracking (s) & 7 & 7 & 3 \\
roadmap, 8 workers: search without it (s) & 6 & 12 & 16 \\
\quad roadmap plan without it (s) & 27 & 38 & 41 \\
\bottomrule
\end{tabular}
}{(missing)}
\end{table}

\begin{table}[h]
\centering
\footnotesize
\setlength{\tabcolsep}{4pt}
\caption{The base parked off its planned pose (shift in the base's own frame, then a
turn about its vertical axis). Each cell: the tolerance of the route used, the time
of the roadmap plan, and the independent check of the saved path.}
\label{tab:robarchshift}
\IfFileExists{figures/prm_table_robarch_shift.tex}{\begin{tabular}{@{}>{\raggedright\arraybackslash}p{5.0cm} rrr@{}}
\toprule
base shift & B1 & B3 & B4 \\
\midrule
5 cm forward, $+3^\circ$ & 20 mm, 30 s, pass & 20 mm, 35 s, pass & 10 mm, 42 s, pass \\
5 cm back-right, $-3^\circ$ & 20 mm, 28 s, pass & 10 mm, 33 s, pass & 10 mm, 30 s, pass \\
10 cm left, $+5^\circ$ & 20 mm, 29 s, pass & 20 mm, 32 s, pass & 10 mm, 50 s, pass \\
10 cm back-left, $-5^\circ$ & goal unholdable & 20 mm, 33 s, pass & 10 mm, 34 s, pass \\
\bottomrule
\end{tabular}
}{(missing)}
\end{table}

\begin{table}[h]
\centering
\footnotesize
\setlength{\tabcolsep}{4pt}
\caption{One collision check of each kind, on 252 configurations spread between home
and the goal posture (one process, idle machine), and whether the split agrees with
the full check.}
\label{tab:robarchsplit}
\IfFileExists{figures/prm_table_robarch_split.tex}{\begin{tabular}{@{}>{\raggedright\arraybackslash}p{7.4cm} rrr@{}}
\toprule
bar & B1 & B3 & B4 \\
\midrule
full check, real scene (ms) & 21.1 & 23.3 & 22.0 \\
scene-only check, real scene (ms) & 7.2 & 10.6 & 11.1 \\
robot-side check, real scene (ms) & 20.1 & 18.8 & 18.1 \\
full check, precompute scene (ms) & 18.5 & 17.2 & 15.1 \\
\midrule
test configurations colliding: full check (of 252) & 45 & 84 & 109 \\
\quad robot side & 45 & 51 & 56 \\
\quad scene only & 0 & 47 & 63 \\
verdicts where the split disagrees with the full check & 0 of 252 & 0 of 252 & 0 of 252 \\
\bottomrule
\end{tabular}
}{(missing)}
\end{table}

\begin{table}[h]
\centering
\footnotesize
\setlength{\tabcolsep}{4pt}
\caption{Every saved path at its waypoints and between them
(\code{prm_check_solved.py --holding-only}). ``Between'': the path interpolated
straight in joint space at $0.1^\circ$ per joint, as the controller moves between
waypoints. The holding error is the planner's own measure (the larger of the two
``where the other flange should be'' gaps). Joint reversals: how often any joint
turns back, moves under $0.2^\circ$ ignored. Ranges over the paths of a row.}
\label{tab:robarchdense}
\IfFileExists{figures/prm_table_robarch_dense.tex}{\begin{tabular}{@{}llrrrrrrr@{}}
\toprule
bar & paths & count & \shortstack[r]{waypoints\\(count)} & \shortstack[r]{largest joint\\step (deg)} & \shortstack[r]{holding error at\\waypoints (mm)} & \shortstack[r]{holding error\\between (mm)} & \shortstack[r]{rotation\\between (deg)} & \shortstack[r]{joint\\reversals} \\
\midrule
B1 & roadmap, planned base & 3 & 824--918 & 0.30--0.31 & 0.0010--0.0016 & 0.0014--0.0023 & 0.0001 & 3 \\
B1 & roadmap, shifted bases & 3 & 789--1008 & 0.27--0.43 & 0.0005--0.0015 & 0.0012--0.0023 & 0.0001 & 0--2 \\
B1 & BiRRT & 3 & 295--346 & 0.54--1.49 & 0.0008 & 0.0064--0.0404 & 0.0004--0.0014 & 3--6 \\
\midrule
B3 & roadmap, planned base & 3 & 755--1343 & 0.28--0.74 & 0.0007--0.0008 & 0.0016--0.0044 & 0.0000--0.0001 & 0--2 \\
B3 & roadmap, shifted bases & 4 & 846--1786 & 0.25--0.51 & 0.0008--0.0020 & 0.0011--0.0033 & 0.0001 & 0--1 \\
B3 & BiRRT & 3 & 579 & 0.49 & 0.0008 & 0.0018 & 0.0001 & 0 \\
\midrule
B4 & roadmap, planned base & 3 & 1276--1515 & 0.30 & 0.0010--0.0011 & 0.0013--0.0017 & 0.0001 & 0 \\
B4 & roadmap, shifted bases & 4 & 648--2990 & 0.24--0.36 & 0.0010--0.0016 & 0.0014--0.0023 & 0.0001--0.0002 & 0--2 \\
B4 & BiRRT & 3 & 1319--1450 & 1.09--5.01 & 0.0013 & 0.0103--0.0781 & 0.0012--0.0071 & 15--31 \\
\bottomrule
\end{tabular}
}{(missing)}
\end{table}

\subsection{What was measured}

All times below are from runs made one at a time on an otherwise idle machine
(16 logical cores). A first pass ran next to the roadmap precompute and was two to
four times slower; its summaries are kept in
\path{prm/robarch_plans/shared_machine_run/} and are not used here, except where
said.

\paragraph{The roadmap precompute (T1).} A 6000-node roadmap takes 8--11 minutes to
build per bar on eight workers and 26 minutes (B1, 4094 edges) to 114 minutes
(B4, 12\,849 edges) to track; B3's and B4's tracking ran next to other jobs, so their
times in Table~\ref{tab:robarchpre} are upper bounds (at B1's rate B4 would take
about 80 minutes). All 13 goal poses can be held for B3 and B4 and 11 of 13 for B1;
of the 18 carry poses, 18 (B3), 17 (B4) and 10 (B1) can be held. B1 is the long bar
grasped far apart, and the family was found on other bars; ten poses were enough
here, so the per-bar carry search that was planned as a fallback was not needed.
The free graph is in a few hundred pieces with roughly a third of the goal nodes
in the largest one, as in Section~\ref{sec:data}.

\paragraph{The collision split holds, and the scene part is the cheap part (T4).}
On 252 configurations per bar the split agreed with the full check every time, in
both senses of Section~\ref{sec:split} (Table~\ref{tab:robarchsplit}). A scene-only
check costs 7--11\,ms against 21--23\,ms for the full check: the pairs it skips
(robot against itself, against its tools and against its load) are most of the
work. In B4's scene 63 of the 252 configurations touch the built structure and 56
collide on the robot's own side, so neither part is negligible there; in B1's
scene nothing touches the scene at all.

\paragraph{Every bar gets a path, and every path passes the independent check (T2).}
Table~\ref{tab:robarchplans} lists the roadmap planner's runs. In each scene nearly
all carry-pose nodes are free (the carry poses are held close to the robot, away
from the structure), a route within 10\,mm (B4) or 20\,mm (B1, B3) exists without any
tube round, and the fine check at 2\,mm found nothing the 1\,cm waypoints had
missed. For B4 the scene removed 9 of the 23 edges the search asked about, and the
search went around them; for B1 and B3 the first route was clean. The tolerance
only says which roadmap edges the route used: the returned path holds the bar
exactly (the bar slips at most 0.004\,mm in the right hand, measured on the saved
joints). Of the 30 plans made on the idle machine (three planner settings and four
shifted bases per bar with the roadmap, three BiRRT runs per bar), 29 returned a
path, and all 29 passed \code{prm_check_solved.py}; the thirtieth is a goal that
cannot be held (below).

\paragraph{Against the BiRRT (T2).} Table~\ref{tab:robarchcompare}. The three bars
are three different cases for the BiRRT, and its log says which:
\begin{description}[nosep,leftmargin=1.5em]
  \item[B1] The BiRRT's start derivation tracks backward from the goal toward its
        home anchors; for B1 it reached none of them in its three 40\,s shares and
        returned a \emph{partial} start 0.41\,m from the goal, leaving the rest to
        another movement. Its path is therefore 0.41--0.51\,m of bar travel, after
        183--188\,s. The roadmap planner returns a whole path from a carry pose
        (1.44\,m of bar travel) in 89\,s, of which 28\,s is planning.
  \item[B3] The straight track from the horizontal home anchor to the goal is
        collision-free, so the BiRRT skips its search and returns in 63--65\,s. The
        roadmap planner takes 84\,s (36\,s of planning, 45\,s of smoothing) for a
        path of the same length. Nothing to gain here.
  \item[B4] The BiRRT has to search between the built bars: 397, 504 and
        1624\,s (the last after four failed attempts of five minutes); in the
        first pass, on the busy machine, one of three runs failed outright after
        28 minutes. Its paths are long: 1.5--2.5\,m of bar travel and
        560--780$^\circ$ of joint travel. The roadmap planner's whole call takes
        63\,s and its path has 1.17\,m and 190$^\circ$.
\end{description}

\paragraph{Where the roadmap planner's time goes (T2, T5).} With eight workers the
roadmap plan takes 28--36\,s, and 15--17\,s of that is starting the workers (each
process loads the scene from files; a planner that keeps its workers for a session
would pay it once). The search takes 3--7\,s and the fine check 4--8\,s. Writing the
scene files the first time a structure is seen cost another 14--22\,s in earlier
runs; a new base pose rewrites only the small state file. With one worker
everything runs in the caller's process on the whole 92-body cell: 57--93\,s. The
shortcut smoothing that follows, the same code as after the BiRRT, runs in the
caller's process and takes another 17--57\,s; it is the largest single part of the
call (its time budget is only looked at between shortcut attempts, each of which
re-solves and re-checks the rest of the path).

\paragraph{The night's tracking saved little here (T4).} Without the tracking file
the search tracks its route edges online and takes 6--16\,s instead of 3--7\,s
(Table~\ref{tab:robarchcompare}): the first or second route worked in every case, so
only 6 to 23 edges were ever asked about, and eight workers track those in a few
seconds. Tracking every edge within 15\,mm in advance is the most expensive part of
the precompute (26--114 minutes) and bought at most 13\,s online. It would matter
when many routes have to be tried; on these three bars, building the roadmap
alone would have been enough.

\paragraph{A base parked off its planned pose (T3).} Table~\ref{tab:robarchshift}:
11 of the 12 shifted bases were planned in 28--50\,s with the same roadmaps, and all
11 paths pass the independent check. The goal seen from a shifted base is a new
pose; its postures became new nodes (4 to 18 per plan) joined to the roadmap, and
only the edges at those nodes had to be tracked online. The one failure is B1
with the base 10\,cm back-left and turned $-5^\circ$: no arm posture can hold the
bar at the goal from there (\code{prm_goal_unholdable}), which no planner can fix;
two of B1's 13 precomputed goal poses had the same problem.

\paragraph{Between the waypoints, and what the joints do (added 2026-10-13).}
Table~\ref{tab:robarchdense} looks at the 29 saved paths the way the roadmap page
now draws them. Every path holds the bar at its waypoints to within
0.002\,mm, whichever planner made it; that is the IK's own precision. Between the
waypoints, where the controller moves the joints in straight lines, the roadmap
paths stay within 0.005\,mm: their waypoints are the 2\,mm fine walk, at most
$0.74^\circ$ apart on any joint, and a straight joint move that short bends the
bar's pose by nothing measurable. The BiRRT paths, whose smoothing leaves steps of
up to $5^\circ$, reach 0.078\,mm and $0.007^\circ$ between waypoints (B4, third
run); twenty times what the roadmap paths show, and still a hundredth of the
loosest tolerance any route here used. Two things follow. The
holding constraint does not need the 2\,mm cut: a path with $1^\circ$ joint steps
would hold the bar to a few hundredths of a millimetre, so the fine check's cost
(Section~\ref{sec:online}) is paid for the collision check, not for holding.
And the joint evolution separates the planners more than the holding error does:
on the roadmap paths a joint turns back 0--3 times per path (moves under
$0.2^\circ$ ignored), the BiRRT's B4 paths turn back 15--31 times, with every
joint taking part (its B1 and B3 paths, which it did not have to search for,
turn back 0--6 times). A shortcut smoother removes detours but not the
zig-zag of the tree's branches, which is what the 560--780$^\circ$ of joint
travel on B4 are made of. On a finer band ($0.006^\circ$) the roadmap paths show
15--74 reversals of a few thousandths of a degree, the IK's noise on a path cut
every 2\,mm; harmless, but a joint-space smoother could take them out.

\paragraph{A note on the export.} The design was exported again on 2026-10-05 and
its cell now holds a ground slab (\code{ground_WG0}, top face at $z=-0.016$\,m).
Paths planned since then see it as a stationary body in the scene check; the
paths above predate it and were measured by completing their saved start state
with the slab's state from the current export (\code{prm_check_solved.py} says so
in its record). The lowest point of any of them, 0.031\,m, stays above it.

\subsection{What this says, and what is still open}

\begin{enumerate}[nosep,label=(\roman*)]
  \item \textbf{The scheme works as proposed.} A roadmap that knows only the robot,
        its tools and the bar is reused across base poses and scenes with no
        recomputation; the scene costs a few hundred to a thousand cheap checks per
        plan, spread over the workers.
  \item \textbf{It pays where the BiRRT has to search.} On B4 the roadmap planner's
        whole call is 6--26 times shorter than the BiRRT's and its path has a quarter
        to a third of the joint travel; on B3, where the BiRRT finds a straight
        track, the BiRRT is the quicker one; on B1 the roadmap planner returns a
        whole path where the BiRRT returns part of one.
  \item \textbf{Both planners end on another arm posture than the export's.} The
        script hands the goal over as two flange frames, so either planner picks
        the posture: the roadmap paths end $155$--$205^\circ$ (largest joint) from the
        authored \code{target_configuration}, the BiRRT paths $203$--$298^\circ$. The
        movements after the carry were authored from that configuration and must
        be re-solved, or the goal restricted to the authored posture
        (\code{goal_conf}, which \code{plan_on_roadmap} accepts; the headless script
        does not pass it yet).
  \item \textbf{No planner checked the floor on these paths.} The cell of
        2026-10-02 had no ground body. The lowest point of the bar, the joint
        halves and the tools was at the goal itself on every path from the
        planned base pose (B1's ground joint lands at $z = -0.003$\,m by design).
        Two shifted-base paths of B4 went lower than their goal: to 0.058\,m and to
        0.031\,m, against 0.071\,m at the goal. The export of 2026-10-05 added a
        ground slab, so plans made since then check against it; whether a floor
        rule is still wanted in the planners (a slab only catches the bodies the
        collision check knows) is to be decided.
  \item \textbf{The holding constraint is kept far inside tolerance, between the
        waypoints too.} At most 0.005\,mm between the waypoints of a roadmap path
        and 0.078\,mm of a BiRRT path (Table~\ref{tab:robarchdense}); the 2\,mm cut
        is for the collision check. The roadmap paths are monotone joint by joint,
        the BiRRT's searched paths are not.
  \item \textbf{Start-up and smoothing now dominate.} Keeping the workers (and the
        roadmap) alive across bars would remove 15\,s per plan; the smoothing is
        shared with the BiRRT and would need its own look.
  \item \textbf{The precompute is per bar and per grasp}, and its expensive half
        (tracking) was not needed on these bars. A build alone is 8--11 minutes per
        bar on eight workers; the node budget and the 15\,mm tracking limit have not
        been tuned.
  \item \textbf{Three bars are a small sample.} They are the first three of one
        structure; later bars carry into denser scenes, where more routes will be
        blocked and the tracking may earn its cost.
\end{enumerate}
